\section{Discussion}
\label{sec:discussion}

\subsection{Geometric Alignment Predicts Composability}

The absolute cosine similarity $\abscossim$ between direction pairs is the
strongest predictor of steering composition quality. The relationship is
non-linear and depends on the \emph{sign} of the alignment:

\begin{itemize}[leftmargin=*,itemsep=2pt,topsep=2pt]
    \item \textbf{Near-orthogonal} ($\abscossim < 0.3$): Best steering
    composition. Both directions operate independently. 3 of 10 pairs fall
    in this regime.
    \item \textbf{Parallel} ($\cossim > 0.8$): Mixed results. The sum
    reinforces both components, and composition can succeed if neither
    direction dominates. 4 of 10 pairs.
    \item \textbf{Anti-parallel} ($\cossim < -0.8$): Worst composition.
    Directions cancel, producing a weak or incoherent composed vector. 3
    of 10 pairs.
\end{itemize}

The Spearman correlation between orthogonality ($1 - \abscossim$) and probing
preservation is $\rho = -0.70$ ($p = 0.025$). However, this reflects a
limitation of the probing metric: orthogonal directions sum to a 45$^\circ$
vector that loses some signal from both concepts, while anti-parallel directions
produce a vector that still partially captures variance from each. The steering
metric provides a more faithful picture of practical composability.

\subsection{Why Are Most Directions Non-Orthogonal?}

The high cosine similarity between seemingly unrelated concepts (\eg
past/present and big/small: $\cossim = 0.92$) is unexpected and demands
explanation. We identify three possible sources:

\para{Sentence structure confounds.}
Both concept pairs use contrastive prompts with similar syntactic structure. The
mean-difference extraction may capture sentence-level variance (word order,
punctuation, length) rather than purely conceptual information. More
sophisticated extraction methods---such as regression-based approaches or
contrastive prompt designs that control for surface features---would likely
reduce this confound.

\para{Superposition.}
With $\dmodel = 2560$, the model represents far more than 2560 concepts
simultaneously. Directions necessarily overlap in this overcomplete
regime~\citep{elhage2022superposition}. The high alignment we observe is
consistent with the superposition hypothesis: the model packs many features
into a shared space, and simple extraction methods recover directions that are
mixtures of the target concept and other co-activated features.

\para{Extraction method limitations.}
The mean-difference estimator is the simplest possible direction extraction
method. It does not account for shared variance between concepts. Methods such
as sparse autoencoder (SAE) features~\citep{chalnev2024improving}, causal inner
product-based extraction~\citep{park2023linear}, or CCA-based
decomposition~\citep{guo2025quantifying} would likely produce more orthogonal
and more composable directions.

\subsection{Probing vs.\ Steering: Two Faces of Composition}

Our results highlight a sharp divergence between probing-based and
steering-based composition metrics. Probing preservation is uniformly high
(${\sim}0.95$) and does not distinguish between composable and non-composable
pairs (ANOVA $p = 0.46$). Steering reveals stark failures for aligned pairs
(component preservation as low as 0.12).

This divergence has methodological implications. Studies that evaluate
composition solely via probing~\citep{park2024geometry, marks2023geometry}
may overestimate the practical composability of linear directions. Probing
measures whether information is \emph{present} in a representation; steering
measures whether it \emph{causally influences} model behavior. For applications
like model control, steering-based evaluation is essential.

\subsection{Connection to Prior Work}

Our findings are consistent with several theoretical predictions:

\para{Causal separability $\Rightarrow$ orthogonality $\Rightarrow$ composition.}
\citet{park2023linear} predict that causally separable concepts are
orthogonal and non-interfering under steering. Our hierarchical pairs
(mammal/bird + fruit/vegetable, $\abscossim = 0.11$) confirm this: they are
near-orthogonal and compose well. This aligns with \citet{park2024geometry}'s
proof that hierarchically related concepts occupy orthogonal subspaces.

\para{Entanglement $\Rightarrow$ interference.}
The true/false + sentiment pair ($\abscossim = 0.99$) fails catastrophically,
consistent with \citeauthor{park2023linear}'s~(\citeyear{park2023linear})
finding that non-separable concepts share directions. The anti-parallel
alignment suggests these directions encode opposing poles of a shared
evaluative dimension.

\para{Composition is task-dependent.}
\citet{todd2023function} find that function vector composition succeeds for
separable sub-operations and fails for entangled ones. Our concept-level
results mirror this pattern: concepts that can be independently varied
(hierarchical categories) compose, while concepts that co-vary (sentiment and
truth valence) interfere.

\para{Scale and method matter.}
\citet{tan2024analyzing} show high per-sample variance in individual
steerability. Our component asymmetry results (7 of 10 pairs show $>0.2$
difference between components) suggest that composition amplifies this
variance. Following \citeauthor{stolfo2024improving}'s~(\citeyear{stolfo2024improving})
finding that layer-separated application outperforms single-layer summation,
multi-layer composition strategies may improve our results.

\subsection{Limitations}

\para{Single model.}
All results are from \pythia. Generalizability to other architectures and scales
is unknown. \citet{todd2023function} find that larger models compose better, so
our results may represent a lower bound on composability.

\para{Simple direction extraction.}
The mean-difference estimator is the simplest available method. SAE-based
features~\citep{chalnev2024improving} or causal inner product-based
directions~\citep{park2023linear} would likely produce cleaner, more orthogonal
directions.

\para{Small concept set.}
With 10 concepts and 10 pairs, our statistical power is limited (the ANOVA is
non-significant at $n = 3$--$5$ per group). A larger-scale study with 50+
concepts would enable stronger statistical claims.

\para{Limited prompts per concept.}
We use 20 contrastive prompt pairs per concept. The relatively low probe
accuracy for some concepts (mammal/bird: 72.5\%) suggests noisy direction
estimates. More prompts would improve direction quality.

\para{1D probing limitation.}
Projecting onto $\dAB$ and using 1D logistic regression is a weak test. The
composed direction may encode both concepts in a 2D subspace that 1D probing
misses, making our probing preservation metric an underestimate of true
information preservation.

\para{No causal inner product.}
We use standard cosine similarity rather than the causal inner
product~\citep{park2023linear}, which could reveal different orthogonality
patterns by accounting for the model's unembedding structure.
